\section{CPU 与 GPU 共享地址不等于共享可见性}

集成 GPU 常与 CPU 共享物理内存控制器，独立显卡则拥有本地 VRAM，并通过 PCIe 与主存交换数据。两种平台都可能向程序提供统一虚拟地址，但“同一个指针可被两侧解释”只解决地址命名，不能自动解决页驻留、Cache 一致性、访问权限和完成顺序。运行时必须知道当前页面位于主存还是 VRAM、哪一侧持有可写副本，以及在所有权转换前有哪些写仍停留在 Cache、写合并缓冲或 DMA 队列中。

CPU 准备输入缓冲时，先写普通 Cache line，再由驱动或运行时执行平台要求的同步操作，使数据对设备可见；随后提交命令描述符，执行发布屏障，最后写 doorbell。GPU 只有在命令处理器取到稳定描述符和资源表后才能执行。内核完成计算并写结果后，fence 值的更新是 GPU 向主机发布完成的边界；CPU 等待 fence 成功后还要按平台一致性规则取得最新数据，才能读取结果。

若系统提供硬件一致性，探测协议可以使 CPU 与 GPU Cache 对同一 Cache line 保持一致，但 fence 仍然必要：一致性回答单个内存位置最终看到哪个值，fence 回答一批命令和结果何时全部完成。若系统不提供一致性，运行时需要显式 flush、invalidate 或使用不可缓存/一致性映射。把“总线支持一致性”理解为“不再需要同步”，会让 CPU 在 GPU 尚未结束时读取一个完全一致但尚未更新的旧值。

\begin{longtable}{L{0.18\linewidth}L{0.28\linewidth}L{0.34\linewidth}L{0.12\linewidth}}
\toprule
边界 & 必须稳定的状态 & 可见结果 & 典型错误 \\
\midrule
CPU 发布输入 & 数据、描述符、资源地址与权限 & GPU 读取同一版本输入 & doorbell 早于描述符可见 \\
地址转换 & 页表、PASID/进程身份、IOMMU/设备 TLB & 设备地址命中允许的物理页 & 页面已回收但设备仍缓存翻译 \\
GPU 执行 & 命令状态、线程块、VRAM/共享内存 & 结果写入目标缓冲 & 内核越界或设备异常 \\
GPU 发布完成 & 结果写入顺序与 fence 值 & CPU/后续队列可确认完成 & fence 更新早于结果可见 \\
显示扫描 & 帧缓冲地址、格式、stride 与翻页时刻 & 显示控制器扫描完整帧 & 渲染与扫描使用不同帧代际 \\
复位恢复 & 队列 generation、映射与在途 DMA & 旧完成不能污染新任务 & 复位后沿用旧指针和 fence \\
\bottomrule
\end{longtable}

\section{页面迁移、设备缺页与超额使用}

统一内存可以在首次访问时把页面迁到使用者附近。GPU 访问尚未驻留的页时，设备 MMU 产生 fault，暂停相关 wave/warp，把地址和访问类型报告给主机；操作系统或驱动分配目标显存、复制页面、更新页表并失效设备 TLB，随后重放访问。这个过程把普通 load 的延迟从纳秒扩展到微秒甚至更高，因此内核即使计算量不大，也可能因频繁迁移而出现明显尾延迟。

若活跃工作集超过 VRAM，页面会在主存与显存间往返。迁移粒度通常是页，而线程访问粒度可能只是几十字节；离散访问会放大 PCIe 流量。运行时可以预取页面、固定关键缓冲、批量安排相邻工作或让部分数据直接驻留主存，但每种策略都要量化带宽与容量。所谓“统一”减少了显式复制代码，不会消除物理链路和容量差异。

设备缺页还必须与 CPU 页框生命周期协调。CPU 不能在 GPU 仍持有翻译或在途访问时回收页面；驱动需要先停止新提交，等待 fence，失效设备页表/IOMMU 映射并确认失效完成，再释放页框。只更新 CPU 页表而遗漏设备 TLB，会让 GPU 使用旧物理地址写入已经分配给其他进程的页面。

\section{命令、计算、显示是三个完成域}

图形程序经常把“命令已提交”“渲染已完成”和“画面已显示”混为一件事。CPU 把命令写入环并敲 doorbell，只表示 GPU 可以开始读取；GPU fence 完成表示对应写入达到设备定义的可见边界；显示控制器在垂直消隐期采用新帧缓冲并扫描最后一个像素，才表示该帧真正出现在接口输出上。三者使用不同队列和时钟域。

双缓冲用 front buffer 供显示扫描、back buffer 供 GPU 渲染。GPU 完成 back buffer 后，显示引擎只能在规定翻页边界交换角色，否则屏幕上半部可能来自旧帧、下半部来自新帧。三缓冲允许 GPU 在显示等待垂直同步时继续渲染，但增加一帧排队可能提高输入到显示的时延。选择缓冲数量时应同时看吞吐、撕裂和交互延迟，不能只看平均帧率。

故障定位也应按三个域读取证据：命令读取指针不动，先查队列、doorbell 和设备状态；计算 fence 不前进，查内核异常、资源映射和 GPU fault；fence 已完成但画面不更新，查翻页队列、帧缓冲格式、显示时序与链路。沿这个顺序可以避免把显示扫描问题误判为着色器计算错误。

\section*{章末练习}

\begin{longtable}{L{0.18\linewidth}L{0.42\linewidth}L{0.30\linewidth}}
\toprule
任务 & 要完成的产物 & 检查要点 \\
\midrule
统一地址 & 解释统一虚拟地址仍需页驻留和同步的原因 & 地址命名不等于数据可见 \\
发布顺序 & 写出 CPU 准备输入到 GPU fence 的顺序 & 描述符先于 doorbell 可见 \\
一致性 & 比较 Cache 一致性与 fence 的职责 & 一个管值，一个管批次完成 \\
页面迁移 & 写出一次 GPU 缺页从暂停到重放的步骤 & 更新页表后还要失效设备 TLB \\
页框回收 & 说明 CPU 何时才能释放 GPU 使用过的页面 & 等待完成并撤销设备翻译 \\
显示边界 & 区分命令提交、渲染完成和显示完成 & 三个队列/时钟域分别取证 \\
\bottomrule
\end{longtable}

\section*{参考解答与要点}

\begin{longtable}{L{0.18\linewidth}L{0.46\linewidth}L{0.26\linewidth}}
\toprule
任务 & 参考解答 & 必须保留的检查点 \\
\midrule
统一地址 & 两侧可以解释同一地址，但页面可能位于不同介质，Cache 和 TLB 也保存独立状态；仍需迁移、权限和所有权协议。 & 指针相同不代表副本最新 \\
发布顺序 & CPU 写数据与描述符，执行所需同步和发布屏障，再写 doorbell；GPU 完成结果写入后更新 fence，CPU 等待并取得可见结果。 & 接受命令不等于完成计算 \\
一致性 & 一致性维持同一 Cache line 的值关系，fence 确认一批异步工作已经越过完成边界。 & 两者不可互相替代 \\
页面迁移 & 设备报告 fault 并暂停相关工作，主机分配/复制页面、更新设备页表、失效 TLB，然后重放访问。 & fault 可显著增加尾延迟 \\
页框回收 & 阻止新提交，等待相关 fence，撤销并确认设备映射/TLB 失效，确认无在途 DMA 后才能释放。 & 不能只改 CPU 页表 \\
显示边界 & doorbell 表示可取命令，GPU fence 表示渲染结果完成，显示完成还要等翻页被采用并扫描输出。 & fence 完成不保证画面已显示 \\
\bottomrule
\end{longtable}
